\scanpage{91}
- \mspage{IV-1010} -
\uline{que, de sorte que les filtrations} $(M'_k)$ \uline{et} $(N'_k)$ \uline{sont les filtrations
L-préadiques. Considérons} $\operatorname{gr}_\bullet(R)$ (resp. $\operatorname{gr}'_\bullet(R)$)
\uline{comme un module gradué sur} $\operatorname{gr}_{\underline K}^\bullet(A)$
(resp. $\operatorname{gr}_{\underline L}^\bullet(A)$).

(i) \uline{Soit} $S$ \uline{une partie de} $\operatorname{gr}_\bullet(R)$
\uline{qui engendre} $\operatorname{gr}_\bullet(R)$ \uline{comme}
$\operatorname{gr}_{\underline K}^\bullet(A)$-\uline{module. Alors la condition c) de}
(19.6.3) \uline{équivaut à la condition suivante} :

c') \uline{L'image de} $S$ \uline{dans} \deleted{\unsure{le}}
$\operatorname{gr}'_\bullet(R)$ \uline{par} $\psi'_R$ \uline{engendre}
$\operatorname{gr}'_\bullet(R)$ \uline{comme} $\operatorname{gr}_{\underline L}^\bullet(A)$-\uline{module}.

(ii) \uline{Supposons vérifiée la condition e) de} (19.6.3), \uline{et de plus que
les quotients des} $\operatorname{gr}_k(R)$ \uline{soient séparés pour la topologie}
$\underline J$-\uline{préadique} (ce qui sera le cas lorsque $A$ est noethérien,
les $f_i$ dans le radical de $A$ et les $\operatorname{gr}_k(N)$ des $A$-modules de type fini).
\uline{Alors, pour
qu'une partie} $S$ \uline{de} $\operatorname{gr}_\bullet(R)$
\uline{formée d'éléments homogènes engendre} $\operatorname{gr}_\bullet(R)$
\uline{comme} $\operatorname{gr}_{\underline K}^\bullet(A)$-\uline{module,
il faut et il suffit que son image par} $\psi'_R$ \uline{engendre}
$\operatorname{gr}'_\bullet(R)$ \uline{comme}
\deleted{$\unsure{\operatorname{gr}_{\underline K}^\bullet(A)}$}
$\operatorname{gr}_{\underline L}^\bullet(A)$-\uline{module}.

(i) Considérons l'homomorphisme \deleted{\unsure{canonique}} de degré 0
\[
\psi_A:(\operatorname{gr}_{\underline K}^\bullet(A)\otimes_A(A/\underline J))[T_1,\ldots,T_n]
 \longrightarrow\operatorname{gr}'_\bullet(A)=\operatorname{gr}_{\underline L}^\bullet(A)
\]
(\uline{0},15.1.5.2). Le fait que cet homomorphisme soit \uline{surjectif} entraîne
que la sous-$\operatorname{gr}_{\underline L}^\bullet(A)$-algèbre de
$\operatorname{gr}'_\bullet(R)$ engendrée par $S'=\psi'_R(S)$ n'est
autre que \deleted{$\unsure{\operatorname{gr}_\bullet(R)}$}
$\operatorname{Im}(\psi'_R)$, d'où l'équivalence de c) et c').

(ii) Puisque e) est vérifiée, il en est de même de c), et en vertu
de (i), il reste à prouver la suffisance de la condition de l'énoncé.
Or, comme $\psi'_R$ est bijectif, dire que $S'$ engendre
$\operatorname{gr}'_\bullet(R)$ considéré comme
$\operatorname{gr}_{\underline L}^\bullet(A)$-module équivaut à dire que $S$ engendre
$(\operatorname{gr}_\bullet(R)\otimes_A(A/\underline J))[T_1,\ldots,T_n]$
comme \deleted{$\unsure{\operatorname{gr}_{\underline K}^\bullet(A)}$}
$\operatorname{gr}_{\underline L}^\bullet(A)$-module, et en vertu
de la surjectivité de $\psi_A$, on peut dans cet énoncé remplacer
l'anneau $\operatorname{gr}_{\underline L}^\bullet(A)$ par
$(\operatorname{gr}_{\underline K}^\bullet(A)\otimes_A(A/\underline J))[T_1,\ldots,T_n]$.
Il revient évidemment
au même de dire que $S$ engendre
$\unsure{\{|f}\operatorname{gr}_\bullet(R)\otimes_A(A/\underline J)$ en tant que module sur
l'anneau $\operatorname{gr}_{\underline K}^\bullet(A)\otimes_A(A/\underline J)$,
ou aussi en tant que $\operatorname{gr}_{\underline K}^\bullet(A)$-module..
\deleted{\unsure{nous}}\inserted{Or,}
\mspage{90}
